\RequirePackage{pdfmanagement}
\SetKeys[document/metadata]{pdfversion=1.7, pdfstandard=A-2u, lang=en-GB}
\input{preamble/pdfA-compliance/glyphtounicode}
\documentclass[11pt,a4paper]{report}

\newcommand*{\PaperTitle}{The Kirchhoff \v{S}olt\'es problem: non-existence theorems and exhaustive searches}
\newcommand*{\PaperTitlePlain}{The Kirchhoff Soltes problem: non-existence theorems and exhaustive searches}
\newcommand*{\PaperShortTitle}{Kirchhoff \v{S}olt\'es problem}
\newcommand*{\PaperSubtitle}{chemistry-open-problems \textperiodcentered{} 24 October 2026 \textperiodcentered{} preliminary (draft by the solver)}
\newcommand*{\PaperCredits}{Solver: Claude Opus 5.5 (effort xhigh) \textperiodcentered{} Verification: pending}
\newcommand*{\PaperKeywords}{Kirchhoff index, resistance distance, \v{S}olt\'es problem, distance-regular graphs, vertex-transitive graphs}
\input{chemop-preamble}
\addbibresource{references.bib}

\newcommand{\Kf}{\operatorname{Kf}}
\newcommand{\Rt}{\operatorname{Rt}}
\newcommand{\one}{\mathbf{1}}

\begin{document}
\makepapertitle

\begin{paperabstract}
Gottwald, Majstorovi\'c Ergoti\'c and Do\v{s}li\'c asked for all graphs $G$ with $\Kf(G-v)=\Kf(G)$ for every
vertex $v$ (Kirchhoff index $=$ sum of effective resistances); only the cycle $C_5$ is known. We prove
non-existence for several classes and extend the published searches. Along a path of degree-two vertices,
$\Kf(G-v_i)$ is a strictly convex quadratic in the position $i$; hence a solution other than $C_5$ with
cyclomatic number $c$ has at most $8(c-1)$ vertices, and a kernel enumeration shows that there is none with
$c\le 5$. Every solution satisfies $a(G)^2\le 2n/\sum_v d_v^{-1}$ ($a$ the algebraic connectivity). Complete
multipartite graphs, strongly regular graphs other than $C_5$, triangle-free distance-regular graphs other
than $C_5$, distance-regular graphs of valency at most four other than $C_5$, all hypercubes and all Hamming
graphs are not solutions; for triangle-free distance-regular graphs we give the closed form
$\Kf(G)-\Kf(G-v)=[(k-1)(n-1)^2-2k(n-k-1)\Kf(G)/n]/[k((k-2)n+2)]$. An exact modular criterion shows that $C_5$
is the only solution on at most $11$ vertices ($910\,914\,360$ two-connected graphs) and that there is no
vertex-transitive solution on $n\le39$ or $n\in\{41,43,45,47\}$ vertices (Holt--Royle census). The problem remains open.
\end{paperabstract}

\section*{The question}
Source \autocite{GMD26}, p.~5 of arXiv:2609.13751:
\begin{openquote}
``Problem 3. Find all graphs $G$ so that the equality $Kf(G) = Kf(G - v)$ holds for all $v \in V(G)$.''
\end{openquote}
and p.~9: ``Our numerical studies failed to identify any additional Kirchhoff \v{S}olt\'es graphs. We
investigated connected graphs with up to 9 vertices. [\dots] we examined all vertex-transitive graphs with up to
23 vertices [\dots]. None of these graphs were Kirchhoff \v{S}olt\'es.'' The problem is the resistance-distance
analogue of \v{S}olt\'es's problem for the Wiener index \autocite{Soltes91,KMS18a,KMS18b,BJM21,AOVVVY23,BKS23}; a
relaxed version is studied in \autocite{MEG25}.

\paragraph{Conventions.} $\Kf(G)=\sum_{u<w}r_G(u,w)=n\sum_{i=1}^{n-1}1/\mu_i$ \autocite{KR93,GM96}, where
$\mu_1\ge\dots\ge\mu_{n-1}>0$ are the non-zero Laplacian eigenvalues; $a(G)=\mu_{n-1}$; $\Gamma=L(G)^+$;
$\Rt_G(v)=\sum_u r_G(u,v)$; $\tau(G)$ is the number of spanning trees; $\Delta_v(G)=\Kf(G)-\Kf(G-v)$. As in
the source, $v$ is \emph{good} if $G-v$ is connected and $\Delta_v=0$, and a \emph{Kirchhoff \v{S}olt\'es (KS)
graph} is a connected graph on $n\ge3$ vertices all of whose vertices are good. Every KS graph is
$2$-connected.

\section{Tools}
\begin{lemma}[note=grounded formula]\label{lem:grounded}
Let $H$ be connected on $N$ vertices, $g\in V(H)$, $A$ the Laplacian of $H$ with the row and column of $g$
deleted and $M=A^{-1}$. Then $\det A=\tau(H)$ and $\Kf(H)=N\operatorname{tr}M-\one^{\mathsf T}M\one$; in
particular $\tau(H)\Kf(H)\in\mathbb Z$.
\end{lemma}
\begin{proof}
$M$ is the Green function grounded at $g$, so $r(u,g)=M_{uu}$ and $r(u,w)=M_{uu}+M_{ww}-2M_{uw}$ for
$u,w\ne g$. Summing over pairs gives $(N-2)\operatorname{tr}M-(\one^{\mathsf T}M\one-\operatorname{tr}M)
+\operatorname{tr}M$. Integrality: $M=\operatorname{adj}(A)/\det A$.
\end{proof}
We use Rayleigh monotonicity, $r_H(x,y)\le d_H(x,y)$, Kron reduction (eliminating a vertex $v$ of degree $d$
replaces its star by a complete graph on $N(v)$ with conductances $1/d$) and the Woodbury identity.

\section{Threads and bounded cyclomatic number}
\begin{lemma}[note=thread formula]\label{lem:thread}
Let $G$ be connected and $P=xv_1\cdots v_{\ell-1}y$ ($\ell\ge2$) a path whose internal vertices have degree
$2$, $x\ne y$, and $R=G-\{v_1,\dots,v_{\ell-1}\}$ connected with $n_R$ vertices. For $a=i-1$, $b=\ell-1-i$,
\[
\Kf(G-v_i)=\Kf(R)+n_R\Bigl[\tbinom{a+1}{2}+\tbinom{b+1}{2}\Bigr]+a\Rt_R(x)+b\Rt_R(y)+\tfrac{a^3-a}{6}
+\tfrac{b^3-b}{6}+\tfrac{ab(a+b+2)}{2}+ab\,r_R(x,y).
\]
For fixed $s=a+b$ this is a polynomial of degree $\le2$ in $a$ with leading coefficient $n_R-1-r_R(x,y)$.
\end{lemma}
\begin{proof}
$G-v_i$ is $R$ with pendant paths of $a$ vertices at $x$ and $b$ vertices at $y$; pendant paths hang on
bridges, so all resistances involving them are sums. Summing the six types of pairs gives the formula.
With $b=s-a$: $n_R[a(a+1)+b(b+1)]/2$ contributes $n_Ra^2$, $(a^3+b^3)/6$ contributes $(s/2)a^2$, and
$ab=sa-a^2$ contributes $-((s+2)/2+r_R(x,y))a^2$.
\end{proof}
\begin{lemma}\label{lem:rpath}
For connected $R$ and $x\ne y$, $r_R(x,y)\le n_R-1$, with equality iff $R$ is a path with ends $x$ and $y$.
\end{lemma}
\begin{proof}
$r_R(x,y)\le d_R(x,y)\le n_R-1$. If equality holds a shortest $x$--$y$ path $P'$ is Hamiltonian, and any
further edge $p_ip_j$ ($j\ge i+2$) puts a unit resistor in parallel with a segment of resistance $j-i$,
giving $r<n_R-1$.
\end{proof}
\begin{corollary}\label{cor:thread}
Let $G$ be $2$-connected but not a cycle. At most two internal vertices of any thread (maximal path of
degree-two vertices) are good; if a thread has exactly two internal vertices and both are good then
$\Rt_R(x)=\Rt_R(y)$. Hence a KS graph is $C_5$ or has all threads with at most two internal vertices.
\end{corollary}
\begin{proof}
The ends $x\ne y$ of a thread have degree $\ge3$, and $R$ is connected (otherwise $x$ would be a cut vertex);
$R$ is not an $x$--$y$ path. By \zcref{lem:thread,lem:rpath}, $a\mapsto\Kf(G-v_{a+1})$ is a strictly convex
quadratic, which takes the value $\Kf(G)$ at most twice. For $\ell=3$ the two values differ by
$\Rt_R(x)-\Rt_R(y)$. A KS cycle $C_n$ needs $(n^3-n)/12=(n-2)(n-1)n/6$, i.e.\ $n=5$ \autocite{GMD26}.
\end{proof}
\begin{theorem}\label{thm:cyclomatic}
A KS graph $G\ne C_5$ with cyclomatic number $c=m-n+1$ satisfies $c\ge2$ and $n\le 8(c-1)$, and is a
subdivision, with $0$, $1$ or $2$ new vertices per edge, of a loopless multigraph without cut vertex, of
minimum degree $\ge3$, with $n_K\le 2(c-1)$ vertices and $m_K\le3(c-1)$ edges. There is no such graph with
$c\le5$; hence every KS graph other than $C_5$ has $m\ge n+5$.
\end{theorem}
\begin{proof}
The kernel (suppress degree-two vertices) is loopless, has no cut vertex and the same cyclomatic number;
$2m_K\ge3n_K$ and \zcref{cor:thread} give the bounds. The kernels are enumerated with nauty
\autocite{MP14} (\texttt{geng -C} for the underlying simple graph, \texttt{multig -e} for multiplicities) and
every admissible subdivision is tested with the exact criterion of \zcref{sec:search}: $c=2$: 7 graphs (the
theta graphs with path lengths in $\{1,2,3\}$); $c=3$: 4 kernels, $1\,038$ graphs; $c=4$: 17 kernels,
$76\,603$ graphs; $c=5$: 118 kernels, $7\,665\,126$ graphs ($5\le n\le32$). None is KS.
\end{proof}

\section{An expander bound and a dense bound}
For $u\ne w$ let $\varphi=\Gamma(e_u-e_w)$ and $E_v(u,w)=\sum_{x\sim v}(\varphi_x-\varphi_v)^2$.
\begin{lemma}\label{lem:pair}
If $G$ and $H=G-v$ are connected and $a_v=a(H)$, then for $u,w\in V(H)$
\[
E_v(u,w)\le r_H(u,w)-r_G(u,w)\le(1+1/a_v)\,E_v(u,w).
\]
\end{lemma}
\begin{proof}
Kron reduction gives $L_{H'}=L_H+S$ with $S=P_N-\frac1d\one_N\one_N^{\mathsf T}$, the orthogonal projector onto
$Y_N=\{y:\operatorname{supp}y\subseteq N(v),\ \sum y=0\}$; write $S=QQ^{\mathsf T}$. On $\one^\perp$,
$X=L_{H'}^{-1}$ and Woodbury give $r_H-r_G=c^{\mathsf T}(I-K)^{-1}c$, $c=Q^{\mathsf T}Xb$,
$K=Q^{\mathsf T}XQ$. From $L_H\succeq a_vI$ we get $L_H+S\succeq a_vI+S$, which acts as $(1+a_v)$ on $Y_N$, so
$0\preceq K\preceq(1+a_v)^{-1}I$ and $I\preceq(I-K)^{-1}\preceq(1+1/a_v)I$. Finally
$\|c\|^2=\sum_{x\sim v}(\varphi_x-\varphi_v)^2$ because $\varphi$ is harmonic at $v$.
\end{proof}
\begin{theorem}[note=expander bound]\label{thm:expander}
Let $G$ be $2$-connected, $a_*=\min_va(G-v)$ and $\varepsilon_v=\sum_{w\ne v}E_v(v,w)\ge(n-1)/d_v$. Then
\[
\sum_v\Delta_v(G)\ \ge\ (1+1/a_*)\sum_v\varepsilon_v-2\Kf(G)/a_*\ \ge\ \tfrac1{a_*}\Bigl[a(G)(n-1)\sum_v d_v^{-1}-2\Kf(G)\Bigr].
\]
Hence $G$ is not KS if $a(G)(n-1)\sum_vd_v^{-1}>2\Kf(G)$, in particular if $a(G)^2>2n/\sum_vd_v^{-1}$; if $G$ is
vertex-transitive, no vertex is good.
\end{theorem}
\begin{proof}
$\Delta_v=\Rt_G(v)-\sum_{u<w;\,u,w\ne v}(r_{G-v}-r_G)$, bounded with \zcref{lem:pair}. One checks
$\sum_{u<w}E_v(u,w)=n\operatorname{tr}(\Gamma^2L^\star_v)$ with $L^\star_v$ the Laplacian of the star at $v$, and
$\sum_vL^\star_v=2L$, so these sums add up to $2\Kf(G)$. The edge currents at $v$ sum to $1$, which gives
$\varepsilon_v\ge(n-1)/d_v$. Further $1+a_*\ge a(G)$ (Fiedler \autocite{Fiedler73}: extend a Fiedler vector of $G-v$
by $0$) and $\Kf(G)\le n(n-1)/a(G)$.
\end{proof}
\begin{corollary}\label{cor:expander}
Every graph with $a(G)^2>2(n-1)$ is not KS; in particular graphs with minimum degree
$\delta>(n-2+\sqrt{2(n-1)})/2$ (since $a\ge2\delta-n+2$), joins $H_1+H_2$ with $\min(n_1,n_2)>\sqrt{2(n-1)}$,
Johnson graphs, Paley graphs $P(q)$ with $q\ge9$, and $d$-regular graphs with second eigenvalue
$\le2\sqrt{d-1}$ and $d\ge12$.
\end{corollary}
\begin{theorem}[note=dense bound]\label{thm:dense}
If $a(G)>1$ then for every $v$,
$\Delta_v(G)\ge1-\sum_{i=1}^{n-1}\frac{n-\mu_i}{\mu_i(\mu_i-1)}$. The bound is attained by the graphs
$K_s+K_{s+1,1}$ of \autocite{GMD26} (half of whose vertices are good).
\end{theorem}
\begin{proof}
With $\Phi(F,N)=\sum_{\theta\in\operatorname{spec}L(F)}\theta/(N-\theta)$, $\Kf(G)=n-1+\Phi(\bar G,n)$ and
$\Kf(G-v)=n-2+\Phi(\bar G-v,n-1)$. Deleting the edges of $\bar G$ at $v$ lowers every eigenvalue (Weyl), and
$\theta/(n-1-\theta)$ is increasing on $[0,n-1)$; so
$\Delta_v\ge1-\sum_\theta\theta/((n-\theta)(n-1-\theta))$ with $\theta=n-\mu$.
\end{proof}

\section{A local deletion formula}
\begin{proposition}\label{prop:local}
Let $v$ be a non-cut vertex, $N=N(v)$, $d=|N|$, $\Gamma_2=\Gamma^2$, and for $a,b\in N$:
$M_{ab}=\Gamma_{ab}-\Gamma_{av}-\Gamma_{vb}+\Gamma_{vv}$,
$(M^2)_{ab}=(\Gamma_2)_{ab}-(\Gamma_2)_{av}-(\Gamma_2)_{vb}+(\Gamma_2)_{vv}+n(\Gamma_{av}-\Gamma_{vv})(\Gamma_{bv}-\Gamma_{vv})$,
$s_a=n(\Gamma_{vv}-\Gamma_{av})$, $t_a=-n[(\Gamma_2)_{av}-(\Gamma_2)_{vv}]-n^2\Gamma_{vv}(\Gamma_{av}-\Gamma_{vv})$,
\[
Z=\begin{pmatrix}M_{NN}-I&s\\s^{\mathsf T}&n-1+n^2\Gamma_{vv}\end{pmatrix},\quad
W=\begin{pmatrix}M^2_{NN}&t\\t^{\mathsf T}&n^2(\Gamma_2)_{vv}+n^3\Gamma_{vv}^2\end{pmatrix}.
\]
Then $\Kf(G-v)=(n-1)[\Rt_G(v)-1-\operatorname{tr}(Z^{-1}W)]$ with $\Rt_G(v)=\operatorname{tr}\Gamma+n\Gamma_{vv}$.
\end{proposition}
\begin{proof}
With $M=L_v^{-1}$ (grounded at $v$), $L_{G-v}=M^{-1}-P_N$ and
$\Kf(G-v)=(n-1)(\operatorname{tr}Y^{-1}-1)$ for $Y=L_{G-v}+J/(n-1)=M^{-1}+UCU^{\mathsf T}$, $U=[E_N,\one]$,
$C=\operatorname{diag}(-I_d,1/(n-1))$; apply Woodbury and express $M$, $M^2$ through $\Gamma$.
\end{proof}
For a distance-regular graph with intersection array $\{b_0,\dots;c_1,\dots\}$, $\Gamma$ and $\Gamma^2$ are
radial and follow from the layer equations $k_ib_i(\psi_i-\psi_{i+1})=\sum_{l\le i}k_lf_l$ \autocite{Biggs93};
two neighbours of $v$ are at distance $1$ or $2$, so $\operatorname{tr}(Z^{-1}W)$ splits over the eigenvalues of the
local graph (\texttt{scripts/drg\_formula.py}, validated against direct exact computation on 17 graphs).

\section{Complete multipartite and strongly regular graphs}
\begin{theorem}\label{thm:multipartite}
For $G=K_{n_1,\dots,n_k}$ (not a star) and $v$ in a part of size $n_j$, $\Delta_v=\phi(n_j)-S$ with
$\phi(x)=\frac{n+x-2}{n-x}+\frac{x(x-1)}{(n-x)(n-1-x)}$ strictly increasing and
$S=\sum_i\frac{n_i(n_i-1)}{(n-n_i)(n-1-n_i)}$. The good vertices form the union of the parts of one size, and
$G$ is never KS: for $K_{k\times s}$, $\Delta_v=\frac{k(k-1)s^2-2ks+2}{(k-1)s((k-1)s-1)}>0$.
\end{theorem}
\begin{theorem}\label{thm:srg}
$C_5$ is the only strongly regular KS graph. For a triangle-free $\mathrm{SRG}(n,k,0,\mu)$,
\[
\Delta_v=\frac{k^4+3k^3\mu-5k^3+3k^2\mu^2-9k^2\mu+5k^2+k\mu^3-3k\mu^2+5k\mu-k+\mu^3-2\mu^2+\mu}{\mu(k+\mu-2)(k^2+k\mu-k+\mu)}>0
\]
unless $(k,\mu)=(2,1)$; for Moore graphs of diameter two $\Delta_v=k(k-2)(k+1)/(k^2+1)$.
\end{theorem}
\begin{proof}
Imprimitive SRGs are complete multipartite (\zcref{thm:multipartite}). Conference graphs with $n\ge9$ have
$a^2=((n-\sqrt n)/2)^2>2k$, and primitive SRGs with integral eigenvalues $r\ge1$, $s\le-2$ have
$\mu=k+rs\le k-2r$, hence $a=k-r\ge(k+1)/2$ and $a^2>2k$ for $k\ge6$; both are excluded by
\zcref{thm:expander}. The feasible primitive parameter sets with $k\le5$ are $(5,2,0,1)$, $(9,4,1,2)$,
$(10,3,0,1)$, $(16,5,0,2)$; the closed form (and, for $(9,4,1,2)$, Jensen and chord bounds on the local-graph
eigenvalue sum, $\Delta\in[64/45,8/5]$) settles them.
\end{proof}

\section{Distance-regular graphs}
\begin{theorem}\label{thm:drg}
(i) If $G$ is distance-regular with $a_1=0$, valency $k\ge2$ and diameter $\ge2$, then at every vertex
\[
\Kf(G)-\Kf(G-v)=\frac{(k-1)(n-1)^2-2k(n-k-1)\Kf(G)/n}{k((k-2)n+2)}.
\]
(ii) $C_5$ is the only KS triangle-free distance-regular graph; for valency $\ge5$ no vertex is good.
(iii) $C_5$ is the only KS distance-regular graph of valency at most four.
(iv) No vertex of a hypercube $Q_d$ or of a Hamming graph $H(d,q)$ ($d,q\ge2$) is good.
(v) A distance-regular graph of valency $\ge3$ with $a(G)\ge2(1+94/101)$ is not KS.
\end{theorem}
\begin{proof}
(i) is \zcref{prop:local} with $b_1=k-1$; the $\Gamma^2$ terms cancel. (ii) For $k\ge5$, Markowsky and Koolen
\autocite[Theorem~1]{MK10} give $r_D\le Kr_1$, $K=1+94/101$, hence $\Kf\le K(n-1)^2/k$, below the threshold
$n(k-1)(n-1)^2/(2k(n-k-1))$ because $2K<4\le k-1$. For $k=3,4$ use the classification lists
\autocite{BCN89,BK99} quoted in \autocite[Cases 2--3]{MK10} and exact evaluation; $k=2$ gives
$\Delta(C_n)=-n(n-1)(n-5)/12$. (iii) The same lists, including the four quartic graphs with $a_1=1$.
(iv) $\Delta(Q_d)>0\iff(d-1)(n-1)^2>d(n-d-1)S_1$, $S_1=\sum_i\binom di/i<4n/(d+1)$; for $H(d,q)$ the analogous
closed form is linear in $\Kf$ and $S_q=\sum_i\binom di(q-1)^i/i<2qn/((d+1)(q-1))$ suffices except for finitely
many cases computed exactly. (v) Combine \zcref{thm:expander} with $\Kf\le K(n-1)^2/k$.
\end{proof}
Closed forms valid for all parameters are also obtained for incidence graphs of symmetric designs, of
generalized quadrangles and hexagons of order $(q,q)$ and for crown graphs (\texttt{check\_drg\_symbolic.py}).

\section{Exhaustive searches}\label{sec:search}
\paragraph{Exact criterion.} For $H\in\{G,G-v\}$, $k_p(H)=N\operatorname{tr}A^{-1}-\one^{\mathsf T}A^{-1}\one\bmod p$
(\zcref{lem:grounded}, LDL$^{\mathsf T}$ over $\mathbb F_p$) equals $\tau(H)\Kf(H)\cdot\tau(H)^{-1}$; if
$k_p(G)\ne k_p(G-v)$ then $p\nmid\tau(G-v)\tau(G)\Kf(G)-\tau(G)\tau(G-v)\Kf(G-v)$, so $\Kf(G)\ne\Kf(G-v)$. Graphs
passing every vertex for two primes are re-checked in exact rationals. Implementations:
\texttt{kscheck.c} ($p=1073741789$, $1073741783$) and \texttt{kscheck\_fast.c} ($p=268435399$, $268435367$),
validated against exact rational values (2\,126 Kirchhoff indices).

\begin{theorem}\label{thm:small}
$C_5$ is the only KS graph with at most $11$ vertices (all $910\,914\,360$ two-connected graphs on
$3\le n\le11$ vertices generated by \texttt{geng -C}; counts equal to OEIS A002218).
\end{theorem}
\begin{theorem}\label{thm:vt}
There is no vertex-transitive KS graph on $n$ vertices for $6\le n\le39$ and $n\in\{41,43,45,47\}$ (every
graph of the corresponding files of the Holt--Royle census \autocite{HR20,HRdata}, all valencies; counts per
$n$ equal to OEIS A006799). The files $n=40,42,44,46$ were not processed.
\end{theorem}


\section{Limitations}
No KS graph other than $C_5$ was found; the general problem remains open. The searches are exact but finite;
floating-point scans of larger families (lattice quotients of $\mathbb Z^2$ up to $1000$ vertices,
generalized Petersen graphs up to $GP(1000,k)$) are explorations only and are not used in any statement.
The valency-$3$ and $4$ classification of distance-regular graphs is used as quoted in \autocite{MK10}.

\paperreferences
\end{document}
